\section{Fermi and Coulomb holes; positivity of $K$ and $J\ge K$}\label{sec:holes}

\begin{reviewclaim}
``\dots two electrons of the same spin cannot occupy the same space \dots\ The corresponding shortage
of the electron of a certain spin in the vicinity of another electron of the same spin is known as the
Fermi hole. Two opposite-spin electrons, on the other hand, are not subject to the exclusion principle.'' (p.~442)
\end{reviewclaim}

\subsection{Pair densities}
The spin-integrated pair density is $P(\mathbf r_1,\mathbf r_2)=N(N-1)\sum_{\sigma_1\sigma_2}\int|\Psi|^2\dd x_3\cdots$.
For two electrons it is $2\sum_{\sigma_1\sigma_2}|\Psi(x_1,x_2)|^2$. The spin parts are normalised, so
\begin{equation}
  P_T(\mathbf r_1,\mathbf r_2)=2\,|\Phi_T|^2=\bigl[\pa(1)\pb(2)-\pb(1)\pa(2)\bigr]^2,\qquad
  P_{\text{OSS}}=\bigl[\pa(1)\pb(2)+\pb(1)\pa(2)\bigr]^2 .
\end{equation}
\begin{enumerate}
  \item Set $\mathbf r_1=\mathbf r_2=\mathbf r$: $P_T(\mathbf r,\mathbf r)=[\pa\pb-\pb\pa]^2=0$ for \emph{every}
        $\mathbf r$. Two same-spin electrons are never found at the same point. This is the Fermi hole,
        and it is exact (no approximation beyond the form of $\Psi$).
  \item $P_{\text{OSS}}(\mathbf r,\mathbf r)=4\,\pa(\mathbf r)^2\pb(\mathbf r)^2\ge0$. It is not zero
        unless the orbitals do not overlap in space. There is no exclusion for an antiparallel pair,
        so the singlet's two electrons can come together. (The Coulomb hole, the reduction of $P$ at
        small $r_{12}$ caused by $1/r_{12}$ itself, is absent from these single-configuration functions.
        It appears only with correlation, i.e.\ configuration mixing.)
\end{enumerate}

\subsection{$K_{ab}\ge0$}
Let $\rho(\mathbf r)=\pa(\mathbf r)\pb(\mathbf r)$ (the overlap density). Then $K_{ab}=\iint\rho(1)\rho(2)/r_{12}$.
\begin{enumerate}
  \item Fourier representation of the Coulomb kernel: $\dfrac1{|\mathbf r|}=\displaystyle\int\frac{\dd^3k}{(2\pi)^3}\frac{4\pi}{k^2}\,e^{i\mathbf k\cdot\mathbf r}$.
  \item Insert this and define $\tilde\rho(\mathbf k)=\int\rho(\mathbf r)e^{-i\mathbf k\cdot\mathbf r}\dd\mathbf r$.
        Because $\rho$ is real, $\tilde\rho(-\mathbf k)=\tilde\rho(\mathbf k)^*$:
        \[
          K_{ab}=\int\frac{\dd^3k}{(2\pi)^3}\frac{4\pi}{k^2}\Bigl(\int\rho(1)e^{i\mathbf k\cdot\mathbf r_1}\Bigr)
          \Bigl(\int\rho(2)e^{-i\mathbf k\cdot\mathbf r_2}\Bigr)
          =\int\frac{\dd^3k}{(2\pi)^3}\frac{4\pi}{k^2}\,|\tilde\rho(\mathbf k)|^2 .
        \]
  \item The integrand is non-negative, so $K_{ab}\ge0$. Equality holds iff $\tilde\rho\equiv0$, i.e.\
        iff $\pa\pb=0$ everywhere (orbitals with no common region of space).
\end{enumerate}

\subsection{$J_{ab}\ge K_{ab}$}
\begin{enumerate}
  \item Take $f(1,2)=\pa(1)\pb(2)-\pb(1)\pa(2)$. Then $\iint f^2/r_{12}\ge0$ because the integrand is non-negative.
  \item Expand: $\iint f^2/r_{12}=J_{ab}+J_{ab}-2K_{ab}=2(J_{ab}-K_{ab})$ (same integrals as in \cref{sec:energies}).
  \item Hence $J_{ab}\ge K_{ab}\ge0$. The triplet repulsion $J-K$ is non-negative, and lies below the
        singlet value $J+K$.
\end{enumerate}
\begin{verifybox}
\texttt{t03\_t05\_hund\_virial\_h2.py} (T3): for all canonical RHF/6-31G orbitals of water,
$\min_{pq}K_{pq}=1.1\times10^{-3}\,E_h>0$ and $J_{pq}\ge K_{pq}$ for every pair. \textbf{PASS}.
\end{verifybox}
